\begin{proposition}[The two nonlocal five-dimensional exceptions]
\label{B:five-nonlocal-exceptions}
Let $N$ be a finite monoid with zero, and let
$A=\mathbb QN/\mathbb Q0$ have dimension five and be nonlocal.
Then one of the following holds:
\begin{enumerate}
\item For all finite-dimensional right and left modules $W,V$,
nonzero $\operatorname{Tor}_2^A(W,V)$ implies nonzero Tor in
arbitrarily high degrees.
\item $A$ is the radical-square-zero algebra of the line quiver
$1\longrightarrow2\longrightarrow3$.
\item $A$ is the two-vertex algebra with
$a\in e_1Ae_2$, $b\in e_2Ae_1$, $ba=0$, and $ab=c\ne0$,
where all products of $c$ with the radical are zero.
\end{enumerate}
In the second case the composite of the two arrows is zero.
In the third case the projective lengths are three and two.
The alternatives are statements about rational algebras; they do
not prescribe a particular multiplicative monoid basis.
\end{proposition}

\begin{proof}
Write $J=\operatorname{rad}(A)$ and $S=A/J$.
Semisimple algebras contribute no positive Tor.
We first dispose of nonsplit semisimple quotients.
Any nonzero $S$-bimodule component incident with a field factor
of degree $d>1$ has rational dimension at least $d$.
The dimension bound therefore leaves only a quadratic field
together with a rational factor and a radical of dimension at
most two, or a larger semisimple quotient with a radical of
dimension at most one.
In the former case a two-dimensional radical component is either
a single cross component, giving a hereditary triangular algebra,
or a loop at the quadratic factor, giving a local square-zero
extension of that field.  If the radical is supported on the
rational factor, that factor has dimension at most three.
The local square-zero syzygy argument of
Lemma~\ref{B:small-quotient-algebra} works over any division
residue algebra, and the other cases are hereditary or products
with already treated small algebras.
A radical of dimension one is supported on rational factors and
gives either a loop or one cross arrow, with the same conclusion.
A matrix or noncommutative division factor of rational dimension
four cannot act nontrivially on a one-dimensional radical, so it
introduces no additional case.

We may now suppose $S\cong\mathbb Q^r$, $r\geq2$.
Lift its primitive orthogonal idempotents.  If $r\geq4$, the
radical has dimension at most one, and the previous argument
applies.  Suppose $r=3$, so $\dim J=2$.
If $J^2\ne0$, the algebra is a product of
$\mathbb Q[x]/(x^3)$ with two rational factors.
If $J^2=0$, it is a radical-square-zero quiver algebra with
three vertices and two arrows.  The proof of
Lemma~\ref{B:nonlocal-four-Tor} applies whenever every vertex
with an entering and a leaving arrow lies on a directed cycle.
The only exception is the directed line on three distinct
vertices, which is alternative~(2).

It remains to treat $r=2$, $\dim J=3$.
If $J^2=0$, every vertex with both an entering and a leaving arrow
lies on a loop or a two-cycle, so the semisimple-syzygy proof of
Lemma~\ref{B:nonlocal-four-Tor} again applies, regardless of
parallel arrows.
If $\dim J^2=2$, the sole arrow in $J/J^2$ is a loop and the
algebra is $\mathbb Q[x]/(x^4)\times\mathbb Q$.
Suppose $\dim J^2=1$.  Nilpotence gives $J^3=0$ and the quiver
has two arrows.

Two parallel cross arrows cannot generate $J^2$.
One loop at each vertex gives
$\mathbb Q[x]/(x^3)\times\mathbb Q[y]/(y^2)$, after possibly
interchanging the vertices.
Two opposite cross arrows give alternative~(3): their two
composites lie at different vertices, so exactly one can span
the one-dimensional space $J^2$.
The remaining configurations are two loops at one vertex, or a
loop and a cross arrow.

Consider the loop-and-arrow case first.  Up to taking the opposite
algebra, write $x\in e_1Je_1$ and $a\in e_1Je_2$.
The only possible products spanning $J^2$ are $x^2$ and $xa$;
they lie in different idempotent components, so precisely one
is nonzero.  Let $P_i=e_iA$ and let $S_i$ be the right vertex
simples.
If $x^2=c\ne0$ and $xa=0$, then
\[
 \operatorname{rad}P_1=L\oplus P_2,
 \quad L=\langle x,c\rangle,\qquad
 \Omega L=S_1\oplus P_2,\quad
 \Omega S_1=L\oplus P_2.
\]
Here $P_2=S_2$ is projective.  Every submodule of a direct sum
of copies of $\operatorname{rad}P_1$ and
$\operatorname{rad}P_2=0$ is a direct sum of $L,S_1,S_2$:
on the first vertex this is the Jordan decomposition for a
square-zero operator, and the cross arrow acts trivially.
Thus every first syzygy has this form.  The displayed identities
make the positive Tor groups of $L$ and $S_1$ two-periodic, so
nonzero degree-two Tor of any pair persists.

If $xa=c\ne0$ and $x^2=0$, then
\[
 \operatorname{rad}P_1=L\oplus P_2,
 \quad L=\langle x,c\rangle,\qquad \Omega L=L.
\]
On $L$, the cross arrow maps its one-dimensional first-vertex
space isomorphically onto its second-vertex space.
Every submodule of $L^s\oplus P_2^t$ therefore has the same
injectivity property and is a direct sum of copies of $L$ and
$P_2$; no $S_1$ summand occurs.  Hence all first syzygies have
this form and positive Tor is one-periodic after one shift.
This also proves persistence.  Taking opposites exchanges the
roles of the two tensor factors and covers the other orientation.

Finally suppose both loops lie at one vertex.  Then
$A\cong B\times\mathbb Q$ with $B$ a four-dimensional local
algebra.  We must use the contracted-monoid hypothesis here:
an arbitrary such local algebra need not satisfy Tor persistence.
We claim that $B$ itself has a contracted monoid basis of
dimension four.

If $N$ has an idempotent $e\notin\{0,1\}$, then $e$ is one of
the two central primitive idempotents of $B\times\mathbb Q$,
because a local algebra has no other idempotents.
If $\dim Ae=4$, the nonzero elements of the submonoid $Ne$
are a basis of $Ae\cong B$; adjoining its zero gives the required
five-element monoid.
If $\dim Ae=1$, then $Ne=\{0,e\}$ is an ideal and the Rees
quotient $N/\{0,e\}$ gives $A/Ae\cong B$ as a contracted
monoid algebra.

Suppose there is no such idempotent.  Every nonunit is nilpotent
to zero.  The finite semigroup of nonunits is nilpotent: otherwise
its descending powers stabilize at a nonzero set, and repeated
factorization in that finite set gives a nonzero element $z$
with $z=zu$ for a product $u$ of nonunits; a sufficiently high
power of the nilpotent element $u$ gives a contradiction.
Thus the nonunit span is a nilpotent ideal and the semisimple
quotient is the rational group algebra of the unit group.
Its dimension is two, so that group is $C_2=\{1,g\}$.
The three other nonzero monoid elements form a basis of $J$.
In the local factor $B$, a unit of order two is $1$ or $-1$:
writing it as $\pm1+j$ reduces the assertion to
$j(\pm2+j)=0$ with $\pm2+j$ invertible.
Multiplication by $g$ permutes the radical monoid basis, so its
action there cannot be $-1$.  It therefore fixes all three
radical basis elements on both sides.
The element $(1+g)/2$ is the identity of $B$, and together with
those three radical basis elements has products either zero or
another member of the same set.  These four elements, together
with zero, again form a finite monoid whose contracted algebra
is $B$.
Theorem~\ref{B:contracted-four-Tor} now proves persistence in
this last case and completes the classification.
\end{proof}
